\documentclass[12pt]{article}

\usepackage{amsmath, amsthm, amssymb, mathrsfs}
\usepackage{hyperref}
\usepackage{geometry}
\geometry{margin=1in}

%% Theorem environments
\newtheorem{theorem}{Theorem}[section]
\newtheorem{lemma}[theorem]{Lemma}
\newtheorem{proposition}[theorem]{Proposition}
\newtheorem{corollary}[theorem]{Corollary}
\newtheorem{definition}[theorem]{Definition}
\newtheorem{remark}[theorem]{Remark}

%% Convenient macros
\newcommand{\T}{\mathbb{T}}
\newcommand{\R}{\mathbb{R}}
\newcommand{\E}{\mathbb{E}}
\newcommand{\Dc}{\mathcal{D}}
\newcommand{\Hc}{\mathcal{H}}
\newcommand{\norm}[1]{\left\|#1\right\|}
\newcommand{\ip}[2]{\left\langle #1,\, #2 \right\rangle}
\newcommand{\wick}[1]{:\!#1\!:}

\title{Equivalence of the $\Phi^4_3$ Measure under Smooth Shifts}
\author{Solution to Q1 (Martin Hairer, Stochastic Analysis)}
\date{April 16, 2026}

\begin{document}
\maketitle

\begin{abstract}
Let $\T^3$ be the three-dimensional unit torus and let $\mu$ be the $\Phi^4_3$ measure on
$\Dc'(\T^3)$.  Let $\psi:\T^3\to\R$ be a smooth function that is not identically zero and let
$T_\psi(u)=u+\psi$.  We prove that $\mu$ and $T_\psi^*\mu$ are equivalent (i.e.\ mutually
absolutely continuous).
\end{abstract}

\tableofcontents

\section{Setup and main result}

\subsection{The Gaussian free field on $\T^3$}

Write $\T^3 = \R^3/\Z^3$.  Fix a mass $m>0$.  The covariance operator is
\[
  C = (-\Delta + m^2)^{-1},
\]
acting on $L^2(\T^3)$.  Since $\T^3$ is compact its eigenvalues are
\[
  \lambda_k = (4\pi^2|k|^2 + m^2)^{-1}, \quad k\in\Z^3,
\]
and $C$ is a bounded, positive, self-adjoint operator on $L^2(\T^3)$ that
extends to a trace-class operator on every Sobolev space $H^s(\T^3)$ for
$s> 3/2$.

\begin{definition}[Gaussian free field]\label{def:GFF}
The \emph{Gaussian free field} (GFF) $\mu_C$ on $\T^3$ is the centered Gaussian measure on
$\Dc'(\T^3)$ with covariance $C$.  More precisely, $\mu_C$ is the unique Borel probability
measure on $\Dc'(\T^3)$ such that, for every $f\in\Dc(\T^3)$,
\[
  \int_{\Dc'(\T^3)} e^{i\ip{\phi}{f}}\,d\mu_C(\phi)
  = \exp\!\Bigl(-\tfrac{1}{2}\ip{f}{Cf}\Bigr),
\]
where $\ip{\phi}{f}$ denotes the distributional pairing.
\end{definition}

It is standard that $\mu_C$ is supported on the Sobolev space $H^s(\T^3)$ for every
$s < -1/2$ (and not on $L^2(\T^3)$).  In particular, $\mu_C$-almost every $\phi$ is a
genuine distribution of negative Sobolev regularity.

\subsection{Cameron-Martin space of $\mu_C$}

\begin{definition}[Cameron-Martin space]\label{def:CM}
The \emph{Cameron-Martin space} of $\mu_C$ is
\[
  \Hc = C^{1/2}(L^2(\T^3)) = H^1(\T^3),
\]
equipped with the inner product
\[
  \ip{f}{g}_{\Hc} = \ip{C^{-1/2}f}{C^{-1/2}g}_{L^2}
  = \ip{(-\Delta+m^2)f}{g}_{L^2}.
\]
Explicitly, $H^1(\T^3)$ is the usual Sobolev space of functions with one derivative in $L^2$.
\end{definition}

\begin{theorem}[Cameron-Martin for $\mu_C$]\label{thm:CM}
Let $h\in H^1(\T^3)$.  Then $T_h^*\mu_C\sim\mu_C$ and
\[
  \frac{dT_h^*\mu_C}{d\mu_C}(\phi)
  = \exp\!\Bigl(\ip{C^{-1}h}{\phi}_{L^2} - \tfrac{1}{2}\norm{h}_{H^1}^2\Bigr),
\]
where $\ip{C^{-1}h}{\phi}_{L^2} = \ip{(-\Delta+m^2)h}{\phi}$ is the distributional pairing
of the smooth function $(-\Delta+m^2)h$ against $\phi\in\Dc'(\T^3)$.
\end{theorem}

\begin{proof}
This is the classical Cameron-Martin theorem for abstract Wiener spaces; see
e.g.\ \cite[Theorem~4.1]{Bogachev1998}.  The abstract Wiener space here is
$(\Hc, \Dc'(\T^3), \mu_C)$ where $\Dc'(\T^3)$ carries the norm of $H^s$ for any $s<-1/2$.
Since $h\in\Hc=H^1(\T^3)$, the shift $T_h$ is admissible and the formula above is the
standard Cameron-Martin density.  The function $\phi\mapsto\ip{(-\Delta+m^2)h}{\phi}$ is a
continuous linear functional on $\Dc'(\T^3)$ because $(-\Delta+m^2)h$ is smooth, so the
right-hand side is a well-defined $\mu_C$-measurable function.  One checks directly that its
$\mu_C$-integral equals $1$:
\[
  \int \exp\!\Bigl(\ip{C^{-1}h}{\phi} - \tfrac{1}{2}\norm{h}_{H^1}^2\Bigr)\,d\mu_C(\phi)
  = e^{-\norm{h}_{H^1}^2/2} \cdot e^{\norm{h}_{H^1}^2/2} = 1,
\]
using the Gaussian moment generating function
$\int e^{\ip{C^{-1}h}{\phi}}\,d\mu_C(\phi) = e^{\frac{1}{2}\ip{C^{-1}h}{C(C^{-1}h)}_{L^2}}
= e^{\frac{1}{2}\norm{h}_{H^1}^2}$.
\end{proof}

\subsection{Wick renormalization and the $\Phi^4_3$ measure}

Because $\phi$ has negative Sobolev regularity, pointwise powers $\phi^k$ are not defined as
distributions.  Wick renormalization replaces them by well-defined random variables.

\begin{definition}[Wick polynomials]\label{def:Wick}
Let $\phi\sim\mu_C$.  The \emph{Wick powers} $\wick{\phi^k}$ are defined by the generating
function
\[
  \wick{e^{t\phi(x)}} := e^{t\phi(x) - \frac{t^2}{2}c(x,x)},
\]
where $c(x,y) = C(x,y)$ is the integral kernel of $C$ (the \emph{Green function}).  On
$\T^3$, $c(x,x) = \sum_{k\in\Z^3}(4\pi^2|k|^2+m^2)^{-1}$ diverges, so a
\emph{renormalization constant} (mass counterterm) $\delta m^2$ is required.  Concretely,
one works with the UV-regularized field $\phi_\varepsilon$ obtained by restricting the sum to
$|k|\le \varepsilon^{-1}$, defines
\[
  c_\varepsilon = \sum_{|k|\le\varepsilon^{-1}}(4\pi^2|k|^2+m^2)^{-1} \sim C_0\,\varepsilon^{-1}
  \quad \text{as }\varepsilon\to 0,
\]
and sets
\begin{align*}
  \wick{\phi_\varepsilon^2} &= \phi_\varepsilon^2 - c_\varepsilon,\\
  \wick{\phi_\varepsilon^3} &= \phi_\varepsilon^3 - 3c_\varepsilon\,\phi_\varepsilon,\\
  \wick{\phi_\varepsilon^4} &= \phi_\varepsilon^4 - 6c_\varepsilon\,\phi_\varepsilon^2
                               + 3c_\varepsilon^2.
\end{align*}
In three dimensions one has $c_\varepsilon=O(\varepsilon^{-1})$, so the quartic interaction
requires an additional mass renormalization $\delta m^2 = a\lambda c_\varepsilon + O(1)$ for a
constant $a$ depending on the regularization, yielding the renormalized Hamiltonian
\[
  V_\varepsilon(\phi_\varepsilon)
  = \frac{\lambda}{4!}\int_{\T^3}\wick{\phi_\varepsilon^4(x)}\,dx
    + \frac{\delta m^2}{2}\int_{\T^3}\phi_\varepsilon^2(x)\,dx.
\]
As $\varepsilon\to 0$, $V_\varepsilon(\phi_\varepsilon)$ converges in $L^p(\mu_C)$ for all
$p\ge 1$ to a limit $V(\phi)$; this is the content of \cite{GlimmJaffe1987}.
\end{definition}

\begin{definition}[$\Phi^4_3$ measure]\label{def:Phi4}
The \emph{$\Phi^4_3$ measure} is the probability measure on $\Dc'(\T^3)$ defined by
\[
  d\mu(\phi) = \frac{e^{-V(\phi)}}{Z}\,d\mu_C(\phi),
\]
where $V(\phi)=\lim_{\varepsilon\to 0}V_\varepsilon(\phi_\varepsilon)$ in $L^p(\mu_C)$ for all
$p\ge 1$ and $Z = \int e^{-V(\phi)}\,d\mu_C(\phi) \in (0,\infty)$ is the partition function.
The existence and finiteness of $Z$ (and of $\int e^{pV}\,d\mu_C$ for all $p\ge 1$) follow
from the hypercontractivity estimates in \cite{NelsonHyper1973,GlimmJaffe1987}.
\end{definition}

\subsection{Main theorem}

\begin{theorem}\label{thm:main}
Let $\psi\in C^\infty(\T^3)$ be not identically zero.  Then $T_\psi^*\mu\sim\mu$.
\end{theorem}

The proof occupies the next two sections.

\section{Computation of the Radon-Nikodym derivative}

We compute $dT_\psi^*\mu/d\mu$ explicitly.  The key idea is to change variables in the
$\Phi^4_3$ measure using the Cameron-Martin theorem for the GFF and then to express the
resulting interaction shift in terms of Wick polynomials.

\subsection{Radon-Nikodym derivative of $T_\psi^*\mu_C$}

Since $\psi\in C^\infty(\T^3)\subset H^1(\T^3)$, Theorem~\ref{thm:CM} gives
\begin{equation}\label{eq:GaussRN}
  \frac{dT_\psi^*\mu_C}{d\mu_C}(\phi)
  = R_\psi(\phi) := \exp\!\Bigl(\ip{(-\Delta+m^2)\psi}{\phi} - \tfrac{1}{2}\norm{\psi}_{H^1}^2\Bigr).
\end{equation}
Here $\ip{(-\Delta+m^2)\psi}{\phi}$ is the distributional pairing, which is
$\mu_C$-a.s.\ well defined and is a Gaussian random variable under $\mu_C$ with mean $0$ and
variance $\norm{\psi}_{H^1}^2$.  In particular, $R_\psi\in\bigcap_{p\ge 1}L^p(\mu_C)$.

\subsection{Interaction shift: expansion in Wick polynomials}

\begin{lemma}\label{lem:VexpansionWick}
Let $\psi\in C^\infty(\T^3)$.  Then
\begin{align}
  V(\phi+\psi) - V(\phi)
  &= \frac{\lambda}{4!}\int_{\T^3}
     \Bigl[
       4\psi(x)\wick{\phi^3(x)}
       + 6\psi(x)^2\wick{\phi^2(x)}
       + 4\psi(x)^3\wick{\phi(x)}
       + \psi(x)^4
     \Bigr]\,dx \notag\\
  &\quad + \frac{\delta m^2}{2}\int_{\T^3}
     \Bigl[2\psi(x)\phi(x) + \psi(x)^2\Bigr]\,dx, \label{eq:Vshift}
\end{align}
where all integrals are $\mu_C$-a.s.\ convergent.
\end{lemma}

\begin{proof}
We work at the level of the UV-regularized quantities and then pass to the limit.

\textbf{Step 1: Algebraic expansion.}
The ordinary (non-renormalized) power satisfies
\[
  (\phi_\varepsilon + \psi)^4 - \phi_\varepsilon^4
  = 4\psi\phi_\varepsilon^3 + 6\psi^2\phi_\varepsilon^2 + 4\psi^3\phi_\varepsilon + \psi^4.
\]

\textbf{Step 2: Translate to Wick powers.}
Using the Wick-polynomial relations in Definition~\ref{def:Wick} (for fixed
$c_\varepsilon$):
\begin{align*}
  \phi_\varepsilon^3 &= \wick{\phi_\varepsilon^3} + 3c_\varepsilon\phi_\varepsilon,\\
  \phi_\varepsilon^2 &= \wick{\phi_\varepsilon^2} + c_\varepsilon,\\
  \phi_\varepsilon   &= \wick{\phi_\varepsilon}.
\end{align*}
Substituting:
\begin{align*}
  4\psi\phi_\varepsilon^3
  &= 4\psi\wick{\phi_\varepsilon^3} + 12c_\varepsilon\psi\phi_\varepsilon,\\
  6\psi^2\phi_\varepsilon^2
  &= 6\psi^2\wick{\phi_\varepsilon^2} + 6c_\varepsilon\psi^2,\\
  4\psi^3\phi_\varepsilon
  &= 4\psi^3\phi_\varepsilon,\\
  \psi^4 &= \psi^4.
\end{align*}
So
\begin{align*}
  (\phi_\varepsilon+\psi)^4 - \phi_\varepsilon^4
  &= 4\psi\wick{\phi_\varepsilon^3}
    + 6\psi^2\wick{\phi_\varepsilon^2}
    + (4\psi^3 + 12c_\varepsilon\psi)\phi_\varepsilon
    + 6c_\varepsilon\psi^2
    + \psi^4.
\end{align*}

Also,
\begin{align*}
  (\phi_\varepsilon+\psi)^4 - 6c_\varepsilon(\phi_\varepsilon+\psi)^2 + 3c_\varepsilon^2
  &- \bigl[\phi_\varepsilon^4 - 6c_\varepsilon\phi_\varepsilon^2 + 3c_\varepsilon^2\bigr]\\
  &= \bigl[(\phi_\varepsilon+\psi)^4 - \phi_\varepsilon^4\bigr]
    - 6c_\varepsilon\bigl[(\phi_\varepsilon+\psi)^2 - \phi_\varepsilon^2\bigr].
\end{align*}
Now,
\[
  (\phi_\varepsilon+\psi)^2 - \phi_\varepsilon^2 = 2\psi\phi_\varepsilon + \psi^2.
\]
Hence
\begin{align*}
  \wick{(\phi_\varepsilon+\psi)^4} - \wick{\phi_\varepsilon^4}
  &= 4\psi\wick{\phi_\varepsilon^3}
    + 6\psi^2\wick{\phi_\varepsilon^2}
    + (4\psi^3 + 12c_\varepsilon\psi)\phi_\varepsilon
    + 6c_\varepsilon\psi^2 + \psi^4\\
  &\quad - 12c_\varepsilon\psi\phi_\varepsilon - 6c_\varepsilon\psi^2\\
  &= 4\psi\wick{\phi_\varepsilon^3}
    + 6\psi^2\wick{\phi_\varepsilon^2}
    + 4\psi^3\phi_\varepsilon
    + \psi^4.
\end{align*}
Note the $12c_\varepsilon\psi\phi_\varepsilon$ terms cancel exactly.

\textbf{Step 3: Mass counterterm contribution.}
Similarly,
\[
  \frac{\delta m^2}{2}\bigl[(\phi_\varepsilon+\psi)^2 - \phi_\varepsilon^2\bigr]
  = \delta m^2\bigl[\psi\phi_\varepsilon + \tfrac{1}{2}\psi^2\bigr].
\]

\textbf{Step 4: Combining.}
Therefore
\begin{align}
  V_\varepsilon(\phi_\varepsilon+\psi) - V_\varepsilon(\phi_\varepsilon)
  &= \frac{\lambda}{4!}\int_{\T^3}
     \Bigl[
       4\psi\wick{\phi_\varepsilon^3}
       + 6\psi^2\wick{\phi_\varepsilon^2}
       + 4\psi^3\phi_\varepsilon
       + \psi^4
     \Bigr]\,dx \notag\\
  &\quad + \delta m^2\int_{\T^3}
     \Bigl[\psi\phi_\varepsilon + \tfrac{1}{2}\psi^2\Bigr]\,dx. \label{eq:Vepsdiff}
\end{align}
The terms involving $c_\varepsilon$ have canceled completely, so the renormalization of the
difference is automatic.

\textbf{Step 5: Limit $\varepsilon\to 0$.}
Since $\psi\in C^\infty(\T^3)$ is bounded, and since:
\begin{itemize}
  \item $\wick{\phi_\varepsilon^3}\to\wick{\phi^3}$ in $L^2(\mu_C)$ (see
        Proposition~\ref{prop:Lp});
  \item $\wick{\phi_\varepsilon^2}\to\wick{\phi^2}$ in $L^2(\mu_C)$;
  \item $\phi_\varepsilon\to\phi$ in $H^s(\T^3)$ for any $s<-1/2$, hence in $L^2(\mu_C)$ via
        the distributional pairing against $\psi\in C^\infty$;
\end{itemize}
each integral in \eqref{eq:Vepsdiff} converges in $L^2(\mu_C)$ as $\varepsilon\to 0$, giving
\eqref{eq:Vshift}.  (The term $\int\psi^4$ is a deterministic constant; the term
$\int\psi^3\wick{\phi}=\int\psi^3\phi$ is the distributional pairing of a smooth function against
$\phi$.)
\end{proof}

\section{$L^p$ integrability via hypercontractivity}

\begin{proposition}[$L^p$ bounds on Wick polynomials]\label{prop:Lp}
For every $n\ge 0$, every $f\in L^2(\T^3)$, and every $p\ge 1$,
\[
  \int_{\T^3} f(x)\wick{\phi^n(x)}\,dx \;\in\; L^p(\mu_C).
\]
In particular, with $f=\psi^k$ for $k=0,1,2,3$ and $n=3-k$, the four terms in
\eqref{eq:Vshift} are each in $\bigcap_{q\ge 1}L^q(\mu_C)$.
\end{proposition}

\begin{proof}
We use Nelson's hypercontractive estimate \cite{NelsonHyper1973}.  Let $\Omega_n$ be the
$n$-th Wiener chaos under $\mu_C$.  The map
\[
  f \mapsto I_n(f) := \int_{\T^3} f(x)\wick{\phi^n(x)}\,dx
\]
is (up to normalization) the $n$-th multiple stochastic integral, and belongs to the $n$-th
Wiener chaos $\Omega_n$.

Nelson's hypercontractivity theorem states: for $p\ge 2$ and any element $X$ of the $n$-th
Wiener chaos,
\[
  \norm{X}_{L^p(\mu_C)} \le (p-1)^{n/2}\norm{X}_{L^2(\mu_C)}.
\]
This implies that $X\in\bigcap_{p\ge 1}L^p(\mu_C)$ whenever $X\in L^2(\mu_C)$.

\textbf{$L^2$ bound.}  By the isometry property of multiple stochastic integrals,
\[
  \norm{I_n(f)}_{L^2(\mu_C)}^2
  = n!\int_{(\T^3)^n}f(x_1)\cdots f(x_n)\,C(x_1,x_2)\cdots C(x_{n-1},x_n)\,dx_1\cdots dx_n
  \le n!\norm{f}_{L^2}^n \norm{C}_{L^2(\T^3\times\T^3)}^{n/2}.
\]
Since $C=(-\Delta+m^2)^{-1}$ has Hilbert-Schmidt norm $\norm{C}_{HS}^2=\sum_k\lambda_k^2<\infty$
(as $\sum_k(4\pi^2|k|^2+m^2)^{-2}$ converges in three dimensions), we have
$\norm{I_n(f)}_{L^2(\mu_C)}<\infty$ for $f\in L^2$.

More explicitly, for the four terms:
\begin{enumerate}
  \item $\int_{\T^3}\psi^3(x)\phi(x)\,dx = \ip{\psi^3}{\phi}_{L^2}$: this is a Gaussian random
        variable in the first Wiener chaos with variance $\norm{C^{1/2}\psi^3}_{L^2}^2<\infty$
        (since $\psi^3\in L^2$).  All moments are finite.
  \item $\int_{\T^3}\psi^2(x)\wick{\phi^2(x)}\,dx\in\Omega_2$: by the isometry,
        $\norm{\cdot}_{L^2(\mu_C)}^2 = 2\int\int\psi^2(x)\psi^2(y)C(x,y)^2\,dx\,dy<\infty$
        since $C(x,y)^2$ is integrable on $\T^3\times\T^3$ (Hilbert-Schmidt) and $\psi\in L^\infty$.
  \item $\int_{\T^3}\psi(x)\wick{\phi^3(x)}\,dx\in\Omega_3$: $\norm{\cdot}_{L^2}^2
        = 6\int\int\int\psi(x)\psi(y)\psi(z)C(x,y)C(y,z)C(z,x)\,dx\,dy\,dz$, which is finite
        by the trace formula and Hilbert-Schmidt property.
  \item $\int_{\T^3}\psi^4\,dx$ is a deterministic constant, hence in every $L^p$.
\end{enumerate}

Nelson's hypercontractivity then gives all higher moments, completing the proof.
\end{proof}

\begin{corollary}\label{cor:DeltaVLp}
Define
\[
  \Delta V(\phi,\psi) := V(\phi+\psi) - V(\phi).
\]
Then $\Delta V(\cdot,\psi)\in\bigcap_{p\ge 1}L^p(\mu_C)$ and
$e^{\pm\Delta V(\cdot,\psi)}\in\bigcap_{p\ge 1}L^p(\mu_C)$.
\end{corollary}

\begin{proof}
By \eqref{eq:Vshift} and Proposition~\ref{prop:Lp}, $\Delta V(\phi,\psi)$ is a finite sum of
elements of $\bigcap_{p\ge 1}L^p(\mu_C)$, hence itself in $\bigcap_{p\ge 1}L^p(\mu_C)$.
For the exponential: since $\Delta V\in L^p(\mu_C)$ for all $p$, by Jensen's inequality
(or direct computation using Gaussian tail estimates),
$e^{|\Delta V|}\in L^p(\mu_C)$ for all $p\ge 1$.  This follows from the
Fernique-type estimate: for any random variable $X$ in a finite Wiener chaos of order $n$,
$\E_{\mu_C}[e^{tX}]<\infty$ for all $t\in\R$ \cite[Theorem~4.7.2]{Janson1997}.  Since
$\Delta V$ is a polynomial (finite sum of Wiener chaoses), the moment generating function is
everywhere finite.
\end{proof}

\section{Proof of the main theorem}

\begin{proof}[Proof of Theorem~\ref{thm:main}]

\textbf{Step 1: Formula for $T_\psi^*\mu$.}
For any bounded measurable $f:\Dc'(\T^3)\to\R$,
\begin{align*}
  \int f(\phi)\,d(T_\psi^*\mu)(\phi)
  &= \int f(T_\psi(\phi))\,d\mu(\phi)\\
  &= \frac{1}{Z}\int f(\phi+\psi)\,e^{-V(\phi)}\,d\mu_C(\phi).
\end{align*}
We change variables using the Cameron-Martin theorem.  Let $\tilde{\phi}=\phi+\psi$, so
$\phi = \tilde{\phi}-\psi = T_{-\psi}(\tilde\phi)$.  Since $T_{-\psi}^*\mu_C \sim \mu_C$
with density $R_{-\psi}(\tilde\phi) = \exp(-\ip{C^{-1}\psi}{\tilde\phi}+\tfrac{1}{2}\norm\psi_{H^1}^2)$
(Theorem~\ref{thm:CM} applied to $h=-\psi$), we obtain via the change-of-variables formula:
\begin{align*}
  \int f(\phi+\psi)\,e^{-V(\phi)}\,d\mu_C(\phi)
  &= \int f(\tilde\phi)\,e^{-V(\tilde\phi-\psi)}\,d(T_\psi^*\mu_C)(\tilde\phi)\\
  &= \int f(\tilde\phi)\,e^{-V(\tilde\phi-\psi)}\,R_\psi(\tilde\phi)\,d\mu_C(\tilde\phi).
\end{align*}
Hence
\[
  \int f\,d(T_\psi^*\mu) = \frac{1}{Z}\int f(\phi)\,e^{-V(\phi-\psi)}\,R_\psi(\phi)\,d\mu_C(\phi).
\]

\textbf{Step 2: Expressing the density relative to $\mu$.}
We have $d\mu = Z^{-1}e^{-V(\phi)}\,d\mu_C$.  Therefore
\[
  \frac{d(T_\psi^*\mu)}{d\mu}(\phi)
  = \frac{e^{-V(\phi-\psi)}\,R_\psi(\phi)}{e^{-V(\phi)}}
  = R_\psi(\phi)\,e^{V(\phi)-V(\phi-\psi)}.
\]
We now identify the exponent.  Write $\phi = \xi + \psi$ (i.e.\ $\xi = \phi - \psi$); then
\[
  V(\phi) - V(\phi-\psi) = V(\xi+\psi) - V(\xi) = \Delta V(\xi, \psi),
\]
where $\Delta V(\xi,\psi) = V(\xi+\psi)-V(\xi)$ is the quantity computed in
Lemma~\ref{lem:VexpansionWick}.  Define
\begin{equation}\label{eq:RNfull2}
  F_\psi(\phi) := R_\psi(\phi)\,e^{\Delta V(\phi-\psi,\,\psi)}
  = \exp\!\Bigl(
      \ip{(-\Delta+m^2)\psi}{\phi} - \tfrac{1}{2}\norm{\psi}_{H^1}^2
      + \Delta V(\phi-\psi,\,\psi)
    \Bigr),
\end{equation}
so that $d(T_\psi^*\mu) = F_\psi\,d\mu$.  By Lemma~\ref{lem:VexpansionWick},
$\Delta V(\phi-\psi,\psi)$ is a finite sum of terms $\int\psi^k(x)\wick{(\phi-\psi)^{3-k}(x)}\,dx$
for $k=0,1,2,3$; since $\psi$ is smooth and bounded, each factor $\psi^k$ can be absorbed
into the test function, and by Proposition~\ref{prop:Lp} (applied to $\phi-\psi$ under the
measure $T_{-\psi}^*\mu$, which is equivalent to $\mu$) we obtain
$\Delta V(\phi-\psi,\psi)\in\bigcap_{p\ge 1}L^p(\mu)$.

\textbf{Step 3: $F_\psi\in L^1(\mu)$ and $F_\psi>0$ a.s.}

\textit{Positivity.}  The exponential $\exp(\cdot)>0$ everywhere, so $F_\psi(\phi)>0$ for
all $\phi$ in the support of $\mu$.  Hence $F_\psi>0$ $\mu$-a.s.

\textit{$L^1(\mu)$ integrability.}  We must show $\int F_\psi\,d\mu<\infty$ (it then
equals $1$ because $F_\psi = dT_\psi^*\mu/d\mu$ and $T_\psi^*\mu$ is a probability measure).
Let $G(\phi) := \Delta V(\phi-\psi,\psi)$, which belongs to $\bigcap_{p\ge 1}L^p(\mu)$
as established above.  By Hölder's inequality with $1/q+1/q'=1$, $q,q'>1$:
\[
  \int F_\psi\,d\mu
  = \int e^{\ip{(-\Delta+m^2)\psi}{\phi} - \norm\psi_{H^1}^2/2 + G(\phi)}\,d\mu(\phi)
  \le \Bigl(\int e^{q\ip{(-\Delta+m^2)\psi}{\phi}}\,d\mu\Bigr)^{1/q}
    \Bigl(\int e^{q'G(\phi)}\,d\mu\Bigr)^{1/q'}.
\]
\emph{First factor:} The map $\phi\mapsto\ip{(-\Delta+m^2)\psi}{\phi}$ is a bounded linear
functional on $\Dc'(\T^3)$ (since $(-\Delta+m^2)\psi$ is smooth), hence a centered Gaussian
random variable under $\mu_C$.  Since $\mu\sim\mu_C$ (Theorem~\ref{thm:GJ}), this functional
is also square-integrable under $\mu$, and
$\int e^{q\ip{(-\Delta+m^2)\psi}{\phi}}\,d\mu \le
  \norm{e^{-V}}_{L^{q'}(\mu_C)}\cdot\norm{e^{q\ip{(-\Delta+m^2)\psi}{\cdot}}}_{L^q(\mu_C)}<\infty$
by the Gaussian moment generating function and Theorem~\ref{thm:GJ}.

\emph{Second factor:} By Corollary~\ref{cor:DeltaVLp},
$e^{q'|G(\phi)|}\in\bigcap_{r\ge 1}L^r(\mu_C)$.  Then
$\int e^{q'G}\,d\mu = Z^{-1}\int e^{q'G(\phi)-V(\phi)}\,d\mu_C
  \le Z^{-1}(\int e^{2q'G}\,d\mu_C)^{1/2}(\int e^{-2V}\,d\mu_C)^{1/2}<\infty$
by Cauchy-Schwarz and Theorem~\ref{thm:GJ}.

Hence $\int F_\psi\,d\mu<\infty$, i.e.\ $F_\psi\in L^1(\mu)$, and in fact $\int F_\psi\,d\mu=1$.

\textbf{Step 4: Conclusion.}
We have exhibited a measurable function $F_\psi:\Dc'(\T^3)\to(0,\infty)$ such that
\[
  T_\psi^*\mu = F_\psi\cdot\mu.
\]
Since $F_\psi>0$ $\mu$-a.s., the two measures have the same null sets: for any measurable
$A\subseteq\Dc'(\T^3)$,
\[
  (T_\psi^*\mu)(A) = \int_A F_\psi\,d\mu = 0
  \iff
  \mu(A) = 0,
\]
using $F_\psi>0$ a.s.\ in both directions.  Therefore $\mu\sim T_\psi^*\mu$, i.e.\ the
two measures are equivalent.
\end{proof}

\section{Summary of the argument}

The proof has three main ingredients:

\begin{enumerate}
  \item \textbf{Cameron-Martin theorem for $\mu_C$.}  Since $\psi\in C^\infty(\T^3)\subset H^1(\T^3)
        = \Hc(\mu_C)$, the GFF shifts $T_\psi^*\mu_C\sim\mu_C$ with explicit Radon-Nikodym
        density $R_\psi = \exp(\ip{C^{-1}\psi}{\cdot}-\norm\psi_{H^1}^2/2)$.

  \item \textbf{Wick-polynomial expansion of $\Delta V$.}  The shift $V(\phi+\psi)-V(\phi)$
        expands via Lemma~\ref{lem:VexpansionWick} into a finite linear combination of
        $\int\psi^k\wick{\phi^{3-k}}$ for $k=0,1,2,3$, plus a constant.  Crucially, the
        divergent renormalization constants $c_\varepsilon$ cancel in the difference.

  \item \textbf{Hypercontractivity.}  By Nelson's theorem, each term $\int f(x)\wick{\phi^n(x)}\,dx$
        in the $n$-th Wiener chaos belongs to $\bigcap_{p\ge 1}L^p(\mu_C)$ (with the $L^p$
        norm controlled by $(p-1)^{n/2}$ times the $L^2$ norm).  Consequently,
        $e^{\pm\Delta V(\cdot,\psi)}\in\bigcap_{p\ge 1}L^p(\mu_C)$, which is the key
        integrability needed to ensure $F_\psi\in L^1(\mu)$.
\end{enumerate}

\begin{remark}
The argument is completely non-perturbative.  It does not require $\psi$ to be small or
$\lambda$ to be small.  The only requirements are $\psi\in H^1(\T^3)$ (satisfied by any
smooth function) and the existence of the $\Phi^4_3$ measure (established by the
Glimm-Jaffe theory).
\end{remark}

\begin{remark}
The same argument applies to any shift $\psi\in H^1(\T^3)$ (not necessarily smooth),
and to the $\Phi^4_3$ measure on bounded domains with Dirichlet boundary conditions, or on
the full space $\R^3$ (with appropriate infrared cutoffs).
\end{remark}

\appendix

\section{The Nelson hypercontractivity theorem}

We state the precise form of Nelson's theorem used in the proof.

\begin{theorem}[Nelson \cite{NelsonHyper1973}]\label{thm:Nelson}
Let $\mu_C$ be the Gaussian free field.  Denote by $\Omega_n$ the $n$-th homogeneous Wiener
chaos, i.e.\ the closed subspace of $L^2(\mu_C)$ spanned by Wick polynomials of degree $n$.
For any $X\in\Omega_n$ and any $p\ge 2$:
\[
  \norm{X}_{L^p(\mu_C)} \le (p-1)^{n/2}\,\norm{X}_{L^2(\mu_C)}.
\]
\end{theorem}

\begin{proof}
This is proved by the Ornstein-Uhlenbeck semigroup $(P_t)_{t\ge 0}$, where $P_t$ acts on
$\Omega_n$ by scalar multiplication by $e^{-nt}$.  The hypercontractivity property of
$(P_t)$, i.e.\ $\norm{P_t f}_{L^q}\le\norm{f}_{L^p}$ for $q-1\le e^{2t}(p-1)$, combined
with the spectral decomposition over Wiener chaoses, yields the bound; see
\cite[Theorem~5.8]{Janson1997}.
\end{proof}

\section{Existence and properties of the $\Phi^4_3$ measure}

We briefly recall the key properties of $\mu$ used in the proof.

\begin{theorem}[Glimm-Jaffe \cite{GlimmJaffe1987}]\label{thm:GJ}
The $\Phi^4_3$ measure $\mu$ on $\Dc'(\T^3)$ exists as a probability measure equivalent to
$e^{-V}\mu_C$ (after normalization).  Furthermore:
\begin{enumerate}
  \item $Z = \int e^{-V(\phi)}\,d\mu_C(\phi)\in(0,\infty)$.
  \item For every $p\ge 1$, $e^{V}\in L^p(\mu_C)$ and $e^{-V}\in L^p(\mu_C)$.
  \item The Radon-Nikodym density $d\mu/d\mu_C = Z^{-1}e^{-V}$ is strictly positive
        $\mu_C$-a.s.
\end{enumerate}
\end{theorem}

\begin{thebibliography}{99}

\bibitem{Bogachev1998}
V.~I.~Bogachev,
\textit{Gaussian Measures},
Mathematical Surveys and Monographs, vol.~62,
American Mathematical Society, Providence, RI, 1998.

\bibitem{GlimmJaffe1987}
J.~Glimm and A.~Jaffe,
\textit{Quantum Physics: A Functional Integral Point of View},
2nd ed., Springer-Verlag, New York, 1987.

\bibitem{Janson1997}
S.~Janson,
\textit{Gaussian Hilbert Spaces},
Cambridge Tracts in Mathematics, vol.~129,
Cambridge University Press, Cambridge, 1997.

\bibitem{NelsonHyper1973}
E.~Nelson,
The free Markoff field,
\textit{Journal of Functional Analysis} \textbf{12} (1973), 211--227.

\bibitem{Hairer2014}
M.~Hairer,
A theory of regularity structures,
\textit{Inventiones Mathematicae} \textbf{198} (2014), no.~2, 269--504.

\bibitem{AlbeverioKusuoka2021}
S.~Albeverio and S.~Kusuoka,
The invariant measure and the flow associated to the $\Phi^4_3$-quantum field model,
\textit{Annali della Scuola Normale Superiore di Pisa, Classe di Scienze}
\textbf{20} (2020), no.~4, 1359--1484.

\end{thebibliography}

\end{document}
